\section{Two-prime coefficient balancing without new variables}

The selected source has 389 rational coordinates, including the input
$X$, and 357 quadratic residuals. Its fixed anchor is
\[
 S_0=3^{2364}=c^6,\qquad c=3^{394}.
\]
The preceding arithmetic argument supplies this parameter. The present
change is an exact rational bijection on the already constructed system;
its zero-set equivalence needs no further arithmetic input.

\subsection{The explicit diagonal map}

Write $x_i$ for an old compiled coordinate and $y_i$ for its replacement.
Set
\[
 x_i=2^{b_i}3^{a_i}y_i.
\]
The input has $a_0=b_0=0$. All unlisted coordinates also have both
exponents zero. The complete nontrivial map is as follows:
\begin{center}
\begin{tabular}{lrrr}
\toprule
Stored source coordinates & Number & $a_i/394$ & $b_i$\\\midrule
The two coefficients of $h_S$ &2&1&0\\
The stored real coefficient of $x(S)$ &1&1&1\\
The stored $\rho$ coefficient of $x(S)$ &1&1&0\\
The two coefficients of $f(S-P)$ &2&1&0\\
The two coefficients of $w=h_S\tau$ &2&2&0\\
Translated anchor witnesses $U,V,W$ &3&4&0\\
Materialized physical anchor $\Theta$ &1&4&0\\
Stored coordinate for the first nontrivial height scale &1&0&$-21$\\\bottomrule
\end{tabular}
\end{center}
The last row denotes precisely the coordinate named
\texttt{heightProduct.scale0.quadratic}. The additional dyadic factor
in the second row targets precisely
\texttt{S.chart.x.0.quadratic}. Names and source identifiers are bound in
the certificate. Semantic height scales and point functions are recovered
through the original value map after decoding these stored coordinates.
In this source the semantic first height scale is
$A=C_0^2=4x_s$, where $x_s$ is the stored coordinate in the last row.
Thus after the change $A=2^{-19}y_s$, and $y_s=2^{19}C_0^2$
need not itself be a rational square. Every detector square-scale and
parity argument applies to the decoded $A=C_0^2$. The transformed
equations retain precisely that value; no square property is imposed
on the new coordinate $y_s$.

There are thirteen changed coordinates: twelve with a ternary factor
and two with a dyadic factor, with one overlap. Every factor is a fixed
positive nonzero rational number, so the displayed map is a bijection
on $\mathbb Q^{389}$ and fixes $X$. Its inverse divides each coordinate
by the same factor. It introduces no new quantified variable or choice
of a root.

The transformed chart equations carry their corresponding powers of
two and three. In particular the coordinates for $h_S$ and $x(S)$ have
different semantic degrees but are scaled as displayed; decoding the
whole transformed system recovers the original identity $x(S)=h_S^2$.

\subsection{Rational row normalization and the zero projection}

For an integer quadratic source residual $R_j(x,S_0)$, let
\[
 \widetilde R_j(y)=R_j((2^{b_i}3^{a_i}y_i)_i,3^{2364}).
\]
Its coefficients are rational because some $b_i$ are negative. Let
$L_j>0$ be the least common multiple of their reduced denominators,
and let $g_j\ne0$ be the signed integer content of $L_j\widetilde R_j$,
with sign chosen to make the first coefficient in canonical monomial
order positive. Put
\[
 d_j=g_j/L_j\in\mathbb Q^*,\qquad
 Q_j=\widetilde R_j/d_j\in\mathbb Z[y_0,\ldots,y_{388}].
\]
Each $Q_j$ is primitive and has degree at most two. The selected quartic is
\[
 P_{389}(y)=\sum_{j=1}^{357}Q_j(y)^2.
\]
Every $d_j$ is nonzero. Consequently, over $\mathbb Q$,
\[
 P_{389}(y)=0
 \ \Longleftrightarrow\ 
 Q_j(y)=0\ \text{for all }j
 \ \Longleftrightarrow\ 
 R_j((2^{b_i}3^{a_i}y_i)_i,3^{2364})=0\ \text{for all }j.
\]
Together with the diagonal bijection, this proves equality of the
existential projections to $X$. Rowwise rescaling preserves the common
zero set; it need not make the new sum of squares one scalar multiple
of the preceding sum of squares.

The rational divisors matter. In the selected system five $d_j$ are
nonintegers, with reduced denominators $2^{21}$ three times,
$2^{42}$ once, and $2^{13}$ once. The certificate records every divisor's
numerator and positive denominator. The checker reconstructs each
substituted row and verifies the exact rational proportionality and
primitive integer normalization, including these five cases.

\subsection{Where the coefficient saving comes from}

The ternary part is the preceding coefficient balancing: the affected
rows have a span at most two in their powers of $c$, so their squares
involve at most $c^4$. With only that map, the large coefficient of the
stored height-scale coordinate in the physical anchor is
$201326592c^2=3\cdot2^{26}c^2$. The extra factor $2^{-21}$ changes it
to $96c^2$. Its squared contribution is then $9216c^4$.

The extra factor two on the stored real coefficient of $x(S)$ balances
the two shifted-point equations after primitive row normalization.
The complete exact expansion, including contributions shared by
different rows, gives
\begin{equation}\label{r8:maximum-coefficient}
 \boxed{\max|\operatorname{coeff}(P_{389})|
      =28902c^4=28902\cdot3^{1576}.}
\end{equation}
This integer has 757 decimal digits and 2513 bits. The quartic has
26,295 nonzero monomials, degree exactly four, content one, and
coefficient one on $X^4$. All 389 coordinates occur. Every coefficient
in the delivered text table is written as an ordinary decimal integer.

For comparison on the preceding anchor alone, the same two dyadic
factors lower its exact maximum from
$1+40532396646334464\cdot3^{1860}$ to $28902\cdot3^{1860}$,
reducing 905 digits to 892 without changing its coordinate count.
The further decrease to 757 digits in
\eqref{r8:maximum-coefficient} comes from the independently justified
anchor exponent 2364.

A linear optimization suggested the factors; its floating-point
output is not a verification dependency or an optimality certificate.
The delivered producer uses the displayed integer exponents and exact
rational arithmetic. A separately authored checker imports no producer
or compiler: it recovers row normalization by pivot ratios and common
denominators, verifies every nonzero divisor, expands the sum of squares
by unordered term pairs, and compares every coefficient in both the
JSON polynomial and the decimal text table. The source gates and
earlier compilation are covered by their separate exact checks.

This establishes the finite transformation and coefficient assertions.
The interpretation of the original source's projection as the integers
retains the global arithmetic qualifications listed in the dependency
inventory.
